\section*{Leaf-language conventions (Appendix A)}

In the following, we introduce basic concepts
related to leaf languages, more details can be found in  \cite{ref11,ref23,ref25,ref24,ref29}.
An NTM $M$ with input alphabet $\Gamma$ is \emph{adequate}, if 
(i) for every input $x \in \Gamma^*$, $M$ does not have an infinite 
computation on input $x$, (ii) the finite set of transition tuples
of $M$ is linearly ordered, and (iii) when terminating $M$ prints a symbol $\alpha(q)$ from a finite alphabet $\Sigma$, where $q$ is 
 the current state of $M$.
For an input $x \in \Gamma^*$, we define the computation tree
by unfolding the configuration
graph of $M$ from the initial configuration. 
By condition (i) and (ii), the computation
tree can be identified with a  
finite ordered tree $T(x) \subseteq \mathbb{N}^*$.
For $u \in T(x)$ let $q(u)$ be the $M$-state of 
the configuration that is associated with the 
tree node $u$.
Then, the leaf string $\leaf(M,x)$
is the string $\alpha(q(v_1)) \cdots \alpha(q(v_k)) \in \Sigma^+$, where 
$v_1,\ldots,v_k$ are all leaves of $T(x)$ listed in lexicographic order.

An adequate NTM $M$ is called  \emph{balanced}, if for every input $x \in \Gamma^*$,
$T(x)$ is a complete binary tree.
With a language $K \subseteq\Sigma^*$ we associate 
the language 
$$
\LEAF(M,K) = \{ x \in  \Gamma^* \mid \leaf(M,x) \in K \}.
$$
Finally, we associate two complexity classes with $K \subseteq\Sigma^*$:
\begin{eqnarray*}
\LEAF(K) &=& \{ \LEAF(M,K) \mid \text{$M$ is an adequate polynomial time
    NTM}\}\\
\bLEAF(K) &=& \{ \LEAF(M,K) \mid \text{$M$ is a balanced polynomial time NTM}\}
\end{eqnarray*}
These classes are closed under polynomial time reductions.
We clearly have $\bLEAF(K)  \subseteq \LEAF(K)$.
The following result was shown in \cite{ref29} by padding computation trees to complete binary trees.

\begin{authorclaim}{Lemma 28}
Assume that $K \subseteq \Sigma^*$ is a language such that $\Sigma$ contains a symbol $1$ with the following property:
if $uv \in K$ for $u,v \in \Sigma^*$ then $u1v \in K$. Then $\LEAF(K)  = \bLEAF(K)$. 
\end{authorclaim}
In particular, we obtain the following lemma:

\begin{authorclaim}{Lemma 29} \label{lemma-JMT}
Let $G$ be a finitely generated group and $\Sigma$ a finite standard generating set for $G$. 
Then $\LEAF(\WP(G,\Sigma))  = \bLEAF(\WP(G,\Sigma))$. 
\end{authorclaim}

Moreover, we have:
\begin{authorclaim}{Lemma 30} \label{lemma-f.g.set-not-important}
Let $G$ be finitely generated group and $\Sigma$, $\Gamma$ finite standard generating sets for $G$. 
Then $\LEAF(\WP(G,\Sigma))  = \LEAF(\WP(G,\Gamma))$. 
\end{authorclaim}

\begin{proof}
Consider a language $L \in \LEAF(\WP(G,\Sigma))$. Thus, there exists an adequate polynomial time
 NTM $M$ such that $L = \LEAF(M,\WP(G,\Sigma))$. We modify $M$ as follows: If $M$ terminates and prints
 the symbol $a \in \Sigma$, it enters a small nondeterministic subcomputation that produces the leaf string $w_a$,
 where $w_a \in \Gamma^*$ is a word that evaluates to the same group element as $a$. Let $M'$ be the resulting
 adequate polynomial time NTM.
 It follows that $\LEAF(M,\WP(G,\Sigma)) =  \LEAF(M',\WP(G,\Gamma))$.
\end{proof}
Lemma 30 allows to omit the standard generating set $\Sigma$ in the notations
$\LEAF(\WP(G,\Sigma))$ and $\bLEAF(\WP(G,\Sigma))$. We will always do that.

For a complexity class $\mathcal C$ we denote by $\forall\mathcal C$ the class of all languages $L$ such that there exists a polynomial $p(n)$ and a language $K\in\mathcal C$ with $L=\{u\mid\forall v\in\{0,1\}^{p(|u|)}:u\#v\in K\}$ (hence, for instance $\forall\mathrm P=\mathrm{coNP}$ and $\forall\mathrm{PSPACE}=\mathrm{PSPACE}$).
